\section{Baselines}
\label{sec:baselines}

\subsection{Setup}

We evaluate two classes of baselines under the runner contract of
\S\ref{sec:benchmark}:
%
\begin{enumerate}
  \item \textbf{Rules-only reference} (\texttt{reference-baseline}): a
        deterministic adapter that emits trivially-formed submissions, run
        on all 250 tasks. It validates that no scorer admits empty answers
        and establishes the deterministic floor.
  \item \textbf{Single-turn LLM} (\texttt{deepseek-chat},
        \texttt{claude-sonnet-4.5}, \texttt{glm-4.5}): the task prompt and
        category-specific output instructions are passed to a frontier
        model in one call; the parsed response is scored as a submission.
        Each model is run under two grounding conditions:
        (i)~\emph{zero-shot}, prompt only; (ii)~\emph{+grounding}, with
        vendored CDISC fixtures (\S\ref{sec:dataset}) inlined into the
        prompt under an 8\,KB-per-request budget.
\end{enumerate}

A retrieval-augmented agent (BM25 over the same fixture subset) is
deferred to v0.1.1 and discussed in \S\ref{sec:limitations}.

\subsection{Decoding and budget}

All LLM baselines use temperature $0$, \texttt{top\_p}$=1$, and a 4096-token
generation cap. Each LLM is evaluated on a stratified subset of 25 tasks
(5 per category, seed \texttt{20260427}); the reference baseline runs on
all 250. Wall-clock latency for the 25-task pass with grounding was
422\,s for DeepSeek-Chat, 519\,s for Claude Sonnet~4.5, and 1325\,s for
GLM-4.5; GLM-4.5 incurred 7~upstream read timeouts on long-prompt tasks
that were scored as zero-output submissions.

\subsection{Results}

Table~\ref{tab:baselines} reports the macro-average score per baseline
under both conditions; per-task submissions and the live HTML
leaderboard are at \texttt{docs/leaderboard/}.

\begin{table}[h]
\centering
\small
\caption{Baseline results on ClinProg-Bench v0.1.0. Per-category columns
are reported for the \emph{+grounding} condition. ``ZS'' denotes the
zero-shot macro for comparison. LLM rows are evaluated on a stratified
25-task subset (5 per category, seed \texttt{20260427}); the reference
row covers all 250 tasks. $\Delta$ = +grounding $-$ zero-shot.}
\label{tab:baselines}
\begin{tabular}{lrrrrrrrr}
\toprule
Baseline & T1 & T2 & T3 & T4 & T5 & Macro & ZS & $\Delta$ \\
\midrule
reference-baseline$^\ddagger$ & --   & --   & --   & --   & --   & --            & 0.05 & --     \\
deepseek-chat         & 0.35 & 0.00 & 0.44 & 0.56 & 0.05 & \textbf{0.28} & 0.27 & $+0.01$ \\
claude-sonnet-4.5     & 0.25 & 0.00 & 0.44 & 0.56 & 0.00 & 0.25          & 0.26 & $-0.01$ \\
glm-4.5               & 0.10 & 0.00 & 0.08 & 0.44 & 0.07 & 0.14$^\dagger$ & 0.22 & $-0.08^\dagger$ \\
\bottomrule
\end{tabular}\\
\smallskip
\raggedright\footnotesize{$^\dagger$ Includes 7 upstream read timeouts on
long-prompt tasks; the operational drop reflects provider-side context
handling under fixture inlining, not raw capability.\\
$^\ddagger$ Reference baseline emits trivially-formed submissions and
does not consume fixtures; only the zero-shot column applies.
Per-category zero-shot scores are
T1\,$=\,0.25$, T2\,$=\,0$, T3\,$=\,0$, T4\,$=\,0$, T5\,$=\,0$.}
\end{table}

\subsection{Discussion}

Four patterns are visible in Table~\ref{tab:baselines}. First, the
category ordering is stable across grounding conditions: T4
(documentation) and T3 (specification extraction) are the easiest, T1
(code generation) is harder, and T2 (log review) and T5 (debug patches)
remain near $0$ even with grounding---these categories require the
specific SAS-log and \texttt{\_v2} program fixtures that the v0.1.0
substrate does not yet author (\S\ref{sec:dataset}). Second, frontier
models cluster within roughly $5$ macro-points of one another in the
zero-shot setting (\(0.22\)--\(0.27\)), suggesting that prompt
engineering and grounding budgets---not raw capability---dominate at
v0.1.0. Third, fixture grounding does \emph{not} uniformly help:
DeepSeek-Chat gains $+0.01$ (T1 lifts from $0.20$ to $0.35$ once spec
documents are inlined), Claude Sonnet~4.5 is unchanged within noise, and
GLM-4.5 loses $0.08$ macro-points operationally because longer prompts
push tail latencies past the provider's read-timeout. We report the
operational number rather than retrying; see \S\ref{sec:limitations} for
the implications and the planned mitigation. Fourth, the identical T4
score (\(0.56\)) across DeepSeek and Claude is a scorer ceiling, not
chance: the documentation rubric awards partial credit for length and
citation structure that any competent LLM satisfies, but the strict
``pass'' threshold requires evidence-grounded claims that none of the
three models produces under either condition. Closing this gap is the
focus of the retrieval-augmented baseline planned for v0.1.1.
